\begin{abstract}
El rápido despliegue de mega-constelaciones de satélites en órbita baja (LEO) para servicios de telecomunicaciones ha incrementado significativamente el riesgo de colisiones orbitales, demandando sistemas avanzados de evasión autónoma y gestión de tráfico espacial (STM). Este artículo presenta una propuesta técnico-metodológica para la evasión de colisiones basada en algoritmos predictivos, sensores embarcados y coordinación multi-agente, orientada a mejorar la seguridad orbital. Se analiza el estado del arte en conjunction assessment, se propone una arquitectura híbrida ground-onboard con aprendizaje por refuerzo para maniobras óptimas, y se evalúan resultados mediante simulaciones de alta fidelidad en escenarios Starlink-like. Los resultados demuestran reducciones sustanciales en probabilidad de colisión y consumo de propelente, contribuyendo a la sostenibilidad de operaciones en LEO. Se discuten implicaciones para estándares IEEE y normativas internacionales de STM.
\end{abstract}